\subsection{Tối ưu hóa nâng cao: Thuật toán chia khối (Tiling/Blocked LU)}
Để giải quyết triệt để các giới hạn vật lý về băng thông bộ nhớ và hiện tượng tràn cache ở các kích thước ma trận cực đại ($N \ge 3000$), ta thực hiện phương án tối ưu nâng cao: \textbf{Khử Gauss phân hoạch chia khối} (Blocked Gaussian Elimination). Nguyên lý thiết kế của giải pháp này tương tự cấu trúc phân rã LU trong thư viện LAPACK chuẩn (\texttt{dgetrf}) phục vụ tính toán hiệu năng cao (HPC).

\subsubsection{Nguyên lý thiết kế và Phân tích Cache}
Ma trận được chia thành các dải cột (Panel) có độ rộng là $B$ (ở đây chọn $B=64$, tức là khối con $64 \times 64$ chiếm khoảng $64 \times 64 \times 8 \text{ bytes} \approx 32\text{ KB}$ dữ liệu số thực \texttt{double}, vừa vặn trọn vẹn trong Cache L1/L2 của từng nhân CPU).
Giải thuật được chia thành 3 pha thực thi chính:
\begin{itemize}
    \item \textbf{Pha 1 (Panel Factorization)}: Khử Gauss tuần tự cục bộ trên dải cột hiện tại (cột từ $kb$ đến $end-1$) để tính toán khối chốt tam giác dưới $L_{21}$ và cập nhật cục bộ dải cột đó. Pha này được thực thi đơn luồng thông qua chỉ thị \texttt{\#pragma omp single}.
    \item \textbf{Pha 2 (Right-Looking Update for Panel Rows)}: Giải thế cho phần khối tam giác trên $U_{12}$ tương ứng với các cột còn lại từ $end$ đến $n-1$. Pha này cũng chạy tuần tự trong khối \texttt{single}.
    \item \textbf{Pha 3 (Trailing Matrix Update - Song song hóa tối đa)}: Cập nhật song song phần ma trận còn lại $A_{22} \leftarrow A_{22} - L_{21} \times U_{12}$ bằng chỉ thị chia sẻ công việc \texttt{\#pragma omp for schedule(static)}. Pha này chiếm trên 90\% tổng khối lượng tính toán của toàn bộ thuật toán. Do các khối con $L_{21}$ và $U_{12}$ đã được nạp sẵn vào bộ nhớ đệm L1/L2 của các lõi ở các pha trước, các luồng có thể liên tục tái sử dụng dữ liệu trực tiếp từ bộ nhớ đệm tốc độ cao, giảm thiểu việc truy xuất bộ nhớ ngoài (RAM).
\end{itemize}

Dưới đây là phần triển khai mã nguồn tối ưu hóa giải thuật khử Gauss chia khối:

\begin{lstlisting}[language=C++, caption={Thuật toán khử Gauss tối ưu hóa bằng phân hoạch chia khối (Blocked LU)}]
vector<double> solveGaussBlocked(const vector<double>& A_flat, const vector<double>& B, int n, int numThreads, int blockSize) {
    vector<double> tempA = A_flat;
    vector<double> tempB = B;
    omp_set_num_threads(numThreads);

    #pragma omp parallel
    {
        for (int kb = 0; kb < n; kb += blockSize) {
            int end = min(kb + blockSize, n);

            // Pha 1 & Pha 2: Tuan tu - Factor panel va cap nhat U12
            #pragma omp single
            {
                for (int i = kb; i < end; ++i) {
                    int pivot = i;
                    for (int j = i + 1; j < n; ++j) {
                        if (abs(tempA[j * n + i]) > abs(tempA[pivot * n + i])) 
                            pivot = j;
                    }
                    if (pivot != i) {
                        for (int k = 0; k < n; ++k) 
                            swap(tempA[i * n + k], tempA[pivot * n + k]);
                        swap(tempB[i], tempB[pivot]);
                    }

                    double pivotVal = tempA[i * n + i];
                    for (int j = i + 1; j < n; ++j) {
                        double factor = tempA[j * n + i] / pivotVal;
                        tempA[j * n + i] = factor;
                        for (int k = i + 1; k < end; ++k) {
                            tempA[j * n + k] -= factor * tempA[i * n + k];
                        }
                        tempB[j] -= factor * tempB[i];
                    }
                }

                for (int k = kb; k < end; ++k) {
                    for (int i = k + 1; i < end; ++i) {
                        double factor = tempA[i * n + k];
                        for (int j = end; j < n; ++j) {
                            tempA[i * n + j] -= factor * tempA[k * n + j];
                        }
                    }
                }
            } // Implicit barrier dong bo hoa truoc khi cap nhat song song

            // Pha 3: Song song - Cap nhat khoi A22 -= L21 * U12
            #pragma omp for schedule(static)
            for (int i = end; i < n; ++i) {
                for (int k = kb; k < end; ++k) {
                    double factor = tempA[i * n + k];
                    for (int j = end; j < n; ++j) {
                        tempA[i * n + j] -= factor * tempA[k * n + j];
                    }
                }
            } // Implicit barrier dong bo hoa truoc khi qua panel tiep theo
        }
    }

    // The nguoc song song hoa reduction
    vector<double> x(n);
    for (int i = n - 1; i >= 0; --i) {
        double sum = 0.0;
        #pragma omp parallel for reduction(+:sum) if(n - i - 1 >= 16)
        for (int j = i + 1; j < n; ++j) {
            sum += tempA[i * n + j] * x[j];
        }
        x[i] = (tempB[i] - sum) / tempA[i * n + i];
    }
    return x;
}
\end{lstlisting}

Tiến hành thực nghiệm so sánh cả 3 phiên bản: Bản gốc (Original), Bản tối ưu phẳng (Optimized) và Bản chia khối (Blocked, $B=64$) trên cấu hình máy Core i5-12400F ở kích thước ma trận lớn nhất $N=5000$. Kết quả đo đạc thực tế được thể hiện trong Bảng~\ref{tab:gauss_blocked_comparison} và trực quan hóa qua Hình~\ref{fig:gauss_blocked_comparison_chart}.

\begin{table}[H]
    \centering
    \caption{Bảng so sánh thời gian thực thi (giây) của 3 thuật toán ($N=5000$)}
    \label{tab:gauss_blocked_comparison}
    \vspace{0.2cm}
    \small
    \begin{tabular}{c|ccc|c}
        \toprule
        \textbf{Threads $p$} & \textbf{Gốc (s)} & \textbf{Tối ưu phẳng (s)} & \textbf{Chia khối (s)} & \textbf{Tốc độ cải thiện cực đại} \\
        \midrule
        $p=1$                & 20.6480          & 20.8525                   & 8.2120                 & 2.51x                             \\
        $p=2$                & 17.9605          & 17.9170                   & 4.4070                 & 4.08x                             \\
        $p=8$                & 15.7935          & 15.6995                   & 2.2880                 & 6.90x                             \\
        $p=12$               & 15.6755          & 15.7165                   & 1.9675                 & 7.97x                             \\
        \bottomrule
    \end{tabular}
\end{table}

\begin{figure}[H]
    \centering
    \begin{tikzpicture}
        \begin{axis} [
                title={So sánh thời gian chạy thực nghiệm: Gốc vs Tối ưu phẳng vs Chia khối ($N=5000$)},
                xlabel={Số lượng luồng $p$},
                ylabel={Thời gian (giây)},
                xtick={1,2,3,4},
                xticklabels={1, 2, 8, 12},
                legend pos=north east,
                ymajorgrids=true,
                grid style=dashed,
                width=0.88\textwidth,
                height=7.0cm,
                ymin=0, ymax=25,
            ]
            \addplot+[mark=square*,thick,red] coordinates {(1, 20.65) (2, 17.96) (3, 15.79) (4, 15.68)};
            \addlegendentry{Bản gốc (Original)}
            \addplot+[mark=triangle*,thick,blue] coordinates {(1, 20.85) (2, 17.92) (3, 15.70) (4, 15.72)};
            \addlegendentry{Tối ưu phẳng (Optimized)}
            \addplot+[mark=otimes*,thick,green!60!black] coordinates {(1, 8.21) (2, 4.41) (3, 2.29) (4, 1.97)};
            \addlegendentry{Chia khối (Blocked, B=64)}
        \end{axis}
    \end{tikzpicture}
    \caption{Biểu đồ thời gian thực thi so sánh thuật toán chia khối tại $N=5000$}
    \label{fig:gauss_blocked_comparison_chart}
\end{figure}

\begin{figure}[H]
    \centering
    \begin{tikzpicture}
        \begin{axis} [
                title={So sánh Speedup: Gốc vs Tối ưu phẳng vs Chia khối ($N=5000$)},
                xlabel={Số lượng luồng $p$},
                ylabel={Độ tăng tốc Speedup $S(p)$},
                xtick={1,2,3,4},
                xticklabels={1, 2, 8, 12},
                legend pos=north west,
                ymajorgrids=true,
                grid style=dashed,
                width=0.88\textwidth,
                height=7.0cm,
                ymin=0, ymax=5,
            ]
            \addplot+[mark=square*,thick,red] coordinates {(1, 1.0) (2, 1.15) (3, 1.31) (4, 1.32)};
            \addlegendentry{Bản gốc (Original)}
            \addplot+[mark=triangle*,thick,blue] coordinates {(1, 1.0) (2, 1.16) (3, 1.33) (4, 1.33)};
            \addlegendentry{Tối ưu phẳng (Optimized)}
            \addplot+[mark=otimes*,thick,green!60!black] coordinates {(1, 1.0) (2, 1.86) (3, 3.59) (4, 4.17)};
            \addlegendentry{Chia khối (Blocked, B=64)}
        \end{axis}
    \end{tikzpicture}
    \caption{Biểu đồ độ tăng tốc Speedup so sánh thuật toán chia khối tại $N=5000$}
    \label{fig:gauss_blocked_speedup_chart}
\end{figure}

\begin{figure}[H]
    \centering
    \begin{tikzpicture}
        \begin{axis} [
                title={So sánh Hiệu suất khai thác luồng: Gốc vs Tối ưu phẳng vs Chia khối ($N=5000$)},
                xlabel={Số lượng luồng $p$},
                ylabel={Hiệu suất khai thác luồng $E(p)$},
                xtick={1,2,3,4},
                xticklabels={1, 2, 8, 12},
                legend pos=north east,
                ymajorgrids=true,
                grid style=dashed,
                width=0.88\textwidth,
                height=7.0cm,
                ymin=0, ymax=1.2,
            ]
            \addplot+[mark=square*,thick,red] coordinates {(1, 1.0) (2, 0.58) (3, 0.16) (4, 0.11)};
            \addlegendentry{Bản gốc (Original)}
            \addplot+[mark=triangle*,thick,blue] coordinates {(1, 1.0) (2, 0.58) (3, 0.17) (4, 0.11)};
            \addlegendentry{Tối ưu phẳng (Optimized)}
            \addplot+[mark=otimes*,thick,green!60!black] coordinates {(1, 1.0) (2, 0.93) (3, 0.45) (4, 0.35)};
            \addlegendentry{Chia khối (Blocked, B=64)}
        \end{axis}
    \end{tikzpicture}
    \caption{Biểu đồ so sánh hiệu suất khai thác luồng tại $N=5000$}
    \label{fig:gauss_blocked_efficiency_chart}
\end{figure}

\textbf{Nhận xét và Đánh giá}:
Phiên bản chia khối (Blocked Gauss) cho thấy sự vượt trội đáng kể về mặt hiệu năng so với hai phiên bản tiền nhiệm:
\begin{itemize}
    \item Ở chế độ đơn luồng ($p=1$), bản chia khối đã nhanh hơn gấp \textbf{2.51 lần} bản gốc nhờ tối ưu hóa cấu trúc lưu trữ và phân bố dữ liệu trong bộ nhớ đệm, giảm thiểu tần suất truy xuất bộ nhớ vật lý.
    \item Khi tăng số lượng luồng, thời gian thực thi giảm mạnh và có xu hướng ổn định. Ở cấu hình 12 luồng, thời gian thực thi rút ngắn xuống còn \textbf{1.9675 giây}, đạt gia tốc nhanh hơn phiên bản gốc \textbf{7.97 lần} (tương đương với mức cải thiện hiệu năng gần 800\%).
    \item Kết quả thực nghiệm này chứng minh: việc tái cấu trúc dữ liệu theo khối (tiling) phù hợp với dung lượng Cache L2/L3 của CPU cho phép giải quyết triệt để nút thắt cổ chai băng thông bộ nhớ RAM, giúp khả năng song song hóa mở rộng tốt theo số lượng luồng.
\end{itemize}

